% Finite chapter. Numerical data remain in the pinned reproduction package.
\section{The finite induction through forty-two}\label{j:finite:chapter}
We prove the finite base by excluding a first violating index. Every
calculation then concerns one attained positive closed-event ratio maximizer
with score $C=C_n>\CE$, and its immediate predecessor has score at most
$\CE$. The computational propositions below assert inequalities on complete
parameter regions; the intervening arguments prove that this same marginal
lies in those regions. No bound on its number of atoms is assumed.

\subsection{Quantitative starting cases}\label{j:finite:starting}

The small cases provide quantitative starting bounds for the later
induction.  The second and third cases use the same idea: a hypothetical
violating extremizer has bounded support, and a polynomial majorizes its
closed sum event on the whole support product.  Independence then evaluates
the polynomial using moments of the same marginal.  No restriction on the
number of atoms is imposed.

\begin{theorem}\label{j:finite:n123}
For every standardized law with finite third absolute moment,
\[
 C_1\le\frac25,\qquad C_2\le\frac{2\sqrt2}{7},\qquad
 C_3\le\frac{59\sqrt3}{250}.
\]
All three upper bounds are strictly smaller than $\CE$.
\end{theorem}

\paragraph{One summand.}
The first numerical assertion is the following checked family of 150 consecutive rational cells $[a,b]$
covering $[-2,1]$, with two rational coefficients $u,v$ in each cell.
It proves
\begin{equation}\label{j:finite:n1-dual}
 L_a-v+uy+vy^2+\frac25|y|^3-\ind_{\{y\le b\}}\ge0
 \quad(y\in\R),\qquad L_a\le\Phi(a).
\end{equation}
Taking expectations cancels the linear and quadratic moment terms.
For $t\in[a,b]$ this gives
\[
 F_P(t)-\Phi(t)\le F_P(b)-L_a\le\frac25\rho.
\]
For $t\le-2$, Cantelli gives $F_P(t)\le1/(1+t^2)\le1/5$.
For $t\ge1$, the discrepancy is at most $1-\Phi(1)<2/5$.
Since $\rho\ge1$, these bounds cover both exterior rays.
Reflection gives the negative branch, including atoms.

\paragraph{A common support argument for two and three summands.}
Let $K_2=2\sqrt2/7$, $K_3=59\sqrt3/250$ and suppose $C_n>K_n$,
where $n=2$ or $3$.  Theorem~\ref{j:core:attainment} supplies a
positive closed-CDF maximizing pair.  Write
\[
 z=H/\sqrt n,\quad F=\Pp(S_n\le H),\quad D_0=F-\Phi(z),
 \quad \lambda=D_0/\rho,\quad \beta=\E[X|X|].
\]
Here $D_0$ denotes the discrepancy, not the support-gap defect used in
the main argument.  The contact condition is equivalently
\begin{align}
 n\Pp(S_{n-1}\le H-y)&\le Q_n(y),\notag\\
 Q_n(y)&=\lambda|y|^3-\frac{z\phi(z)+3D_0}{2}y^2
 -(\sqrt n\phi(z)+3\lambda\beta)y
 +nF+\frac{z\phi(z)+D_0}{2},\label{j:finite:n23-contact}
\end{align}
with equality on the entire closed support.
The published affine estimate
\begin{equation}\label{j:finite:n23-affine}
 \frac{\sqrt n\,\Delta_n(P)}{\rho(P)}\le\frac{1661}{5000}
       \left(1+\frac{429}{1000\rho(P)}\right)
\end{equation}
from \citet{Shevtsova2013} rules out $\rho\ge2$ at $n=2$ and
$\rho\ge15/8$ at $n=3$.  Together with $|\beta|\le1$,
$\phi(z)<2/5$ and $|z\phi(z)|<1/4$, it gives the following lower
bounds $Q_n(y)\ge P_n(|y|)$:
\begingroup\renewcommand{\arraystretch}{1.8}
\[
\begin{array}{c|c|c}
 n&\rho\text{ ceiling}&P_n(r)\\\hline
 2&2&\displaystyle\frac27r^3-\frac{67}{56}r^2-\frac{23}{14}r-\frac18\\[2pt]
 3&15/8&\displaystyle\frac8{35}r^3-\frac45r^2-\frac85r-\frac18
\end{array}
\]
\endgroup
Indeed, for $n=2$, $\lambda>2/7$, $D_0<5/7$ and
$\lambda<5/14$ suffice: $\sqrt2<10/7$ gives
$\sqrt2\phi(z)+3\lambda<4/7+15/14=23/14$.
For $n=3$, the corresponding bounds are
$\lambda>8/35$, $D_0<9/20$ and $\lambda<2/7$.
In the latter case \eqref{j:finite:n23-affine} gives $C_3\rho<77/100$,
and $\sqrt3>12/7$ then gives $D_0=C_3\rho/\sqrt3<9/20$.
The shifted differences are
\begin{align*}
 P_2(6+s)-2&=\frac27s^3+\frac{221}{56}s^2+
                  \frac{104}{7}s+\frac{373}{56},\\
 P_3(6+s)-3&=\frac8{35}s^3+\frac{116}{35}s^2+
                  \frac{472}{35}s+\frac{2197}{280}.
\end{align*}
Every coefficient is positive.  Equality in
\eqref{j:finite:n23-contact} therefore excludes every support point with
$|y|\ge6$, because its left side is at most $n$.

\paragraph{The two-summand majorants.}
Put $\gamma=\E X^3$ and
\[
 v=(1,y,y^2,y^3,|y|^3),\qquad \mu=\E v(X)=(1,0,1,\gamma,\rho).
\]
For a rational symmetric coefficient array indexed by $i\le j$, let
\[
 M_c(x,y)=\sum_i c_{ii}v_i(x)v_i(y)+
 \sum_{i<j}c_{ij}\{v_i(x)v_j(y)+v_j(x)v_i(y)\},
\]
and $q_c(\gamma,\rho)=\sum_i c_{ii}\mu_i^2+
2\sum_{i<j}c_{ij}\mu_i\mu_j$.
Independence gives $\E M_c(X_1,X_2)=q_c(\gamma,\rho)$.
Lemma~\ref{j:core:moments} puts the actual moments in
\[
 \mathcal D_2=\{(\gamma,\rho)\in[-2,2]\times[1,2]:
 \rho^2(\gamma^2+4)-(\gamma^2+2)^2\ge0\}.
\]
The 70 certificate rows use
$a_j=-5/2+j/20$, $b_j=a_j+1/20$, $0\le j<70$.
The checked numerical assertions, in the notation of
Lemma~\ref{j:cert:tensor}, are
\begin{align*}
 M_{c^{(j)}}(x,y)&\ge\ind_{\{x+y\le b_j\}}
                    &&((x,y)\in[-6,6]^2),\\
 q_{c^{(j)}}(\gamma,\rho)&\le\Phi(a_j/\sqrt2)+\frac27\rho
                    &&((\gamma,\rho)\in\mathcal D_2).
\end{align*}
Each sign quadrant is cut by the event line and triangulated into
closed triangles.  Degree-six Bernstein coefficients, with exact
subdivision, establish the first inequality.  The second inequality is
checked on moment rectangles; a rectangle is discarded only if
$\rho_+^2(g_*+4)-(g_*+2)^2<0$, where $g_*$ is the minimum of
$\gamma^2$ on that rectangle.  This is valid because the expression
increases with $\rho$ and decreases with $\gamma^2$ when $\rho\le2$.

For a threshold in $[a_j,b_j]$, expectation of the first inequality and
the second give discrepancy at most $2\rho/7$.
For $H\le-5/2$, Cantelli gives $2/(2+H^2)\le8/33<2/7$;
for $H\ge1$, $1-\Phi(H/\sqrt2)<2/7$.
These bounds contradict the chosen violating maximizer.  This proves
the second assertion of Theorem~\ref{j:finite:n123}.

\paragraph{The three-summand majorants.}
Use the features
\[
 (v_0,\ldots,v_6)(y)=(1,y,y^2-1,y^3,|y|^3,y^4,y^5).
\]
Let $\mathcal I$ be the 59 nondecreasing triples from $\{0,\ldots,6\}$
whose largest entry is less than 5 or which contain an entry 1 or 2.
For $\alpha\in\mathcal I$, average the associated product over its
distinct coordinate permutations, obtaining $V_\alpha$.
Set $M_c=\sum_\alpha c_\alpha V_\alpha$.
The bounded-support argument above makes every feature integrable.
Every term involving a fourth or fifth power contains a zero-mean
factor $X_i$ or $X_i^2-1$ in a different coordinate.  Independence
therefore gives the exact identity
\[
 \E M_c(X_1,X_2,X_3)=q_c(\gamma,\rho)
 =\sum_{\substack{\alpha\in\mathcal I\\\alpha\text{ uses only }0,3,4}}
     c_\alpha\gamma^{\#3(\alpha)}\rho^{\#4(\alpha)}.
\]
There are ten surviving monomials and 49 zero-expectation controls.
The controls improve the pointwise majorant without adding moment
parameters.

The second checked numerical assertion uses 111 rational rows forming a consecutive closed cover of
$[-7/2,3/2]$. They satisfy
\begin{align*}
 M_{c^{(j)}}(x_1,x_2,x_3)&\ge\ind_{\{x_1+x_2+x_3\le b_j\}}
                    &&((x_1,x_2,x_3)\in[-7,7]^3),\\
 q_{c^{(j)}}(\gamma,\rho)&\le\Phi(a_j/\sqrt3)+\frac{59}{250}\rho
                    &&((\gamma,\rho)\in\mathcal D_3),
\end{align*}
where $\mathcal D_3$ replaces $2$ by $15/8$ in the rectangular
bounds defining $\mathcal D_2$.
Symmetry reduces the sign patterns to four ordered-coordinate
orthants.  A box is discarded only when it has no ordered point.
On every other box, degree-five tensor Bernstein bounds establish
the pointwise inequalities.  On a box crossing the event plane,
the included side is checked by $M-1\ge s(x_1+x_2+x_3-b_j)$ with
$s\le0$, and the other side by the analogous inequality with $s\ge0$.
Existence of a slope is checked by intersecting exact rational intervals;
zero plane coefficients are checked directly.  This retains the closed
event boundary.

As in the two-summand case, expectation bounds the discrepancy on each
threshold interval.  The exterior bounds are
$3/(3+H^2)\le12/61<59/250$ for $H\le-7/2$ and
$1-\Phi(H/\sqrt3)<59/250$ for $H\ge3/2$.
They contradict the violating maximizer and prove the third assertion.
The first assertion was proved above.  Rational square-root and Machin
bounds give $K_2<0.405<0.409<\CE$ and $K_3<0.409<\CE$.

\subsubsection{Localization at four summands}
\label{j:finite:n4-localization-section}

The third bound is needed quantitatively.  Merely knowing $C_3\le\CE$
would not give the following support reduction.
Write
\[
 K_3=\frac{59\sqrt3}{250},\quad c_0=\frac{256}{625},\quad
 T=\frac{4097}{10000};\qquad K_3<c_0<T<\CE.
\]
For any standardized law set
$\beta=\E[X|X|]$, $u=\E X_+=\E X_->0$,
$v_\pm=(1\pm\beta)/2$, and
\begin{equation}\label{j:finite:n4-reference}
 a=\frac{v_-}{u},\quad b=\frac{v_+}{u},\quad
 w=\frac{4u^2}{1-\beta^2},\quad
 Q=\frac ua\delta_{-a}+(1-w)\delta_0+\frac ub\delta_b.
\end{equation}
Cauchy's inequality gives $0<w\le1$ and makes $Q$ a standardized
probability law.  Put $r=\rho(Q)=(1+\beta^2)/\sqrt{w(1-\beta^2)}$.
\begin{lemma}\label{j:finite:n4-excess}
The nonnegative quantity
\[
 W(X)=X_+(X_+-b)^2+X_-(X_--a)^2
\]
has expectation $\eta=\rho(P)-r$.  Thus $r\le\rho(P)$, with
equality exactly when $P=Q$.  Moreover, writing $c(X)=b$ for $X>0$
and $c(X)=a$ for $X<0$,
\begin{align}
 \E X^3&=b-a+\E[\operatorname{sgn}(X)W(X)],\notag\\
 \E X^4&=a^2-ab+b^2+\E[(|X|+2c(X))W(X)]
 \label{j:finite:n4-physical-moments}
\end{align}
whenever the fourth moment is finite.
\end{lemma}
\begin{proof}
On either half-line put $y=|X|$ and $c=a$ or $b$.  The identities
$y^3=y(y-c)^2+2cy^2-c^2y$ and
$y^4=(y+2c)y(y-c)^2+3c^2y^2-2c^3y$ integrate using
$\E[y;\text{half-line}]=u$ and
$\E[y^2;\text{half-line}]=uc$.
If $\E W=0$, the support is contained in $\{-a,0,b\}$ and the
first two moments fix the masses.  Conversely $P=Q$ has $W=0$.
\end{proof}
The three-point law is a reference;
the original law remains arbitrary.

\begin{proposition}\label{j:finite:n4-localization}
If $C_4>c_0$, a positive closed-CDF maximizing pair satisfies
\[
 \supp P\subset(-101/25,101/25),\quad\rho(P)<46/25,
 \quad 625/2116<w\le1,\quad|\beta|<321/500.
\]
In particular it belongs to the closed numerical class with
$|X|\le17/4$ and $|\beta|\le13/20$.
\end{proposition}
\begin{proof}
Use Theorem~\ref{j:core:attainment}.  The affine inequality
\eqref{j:external:affine} first gives
$\rho<979013/532500<46/25$, without a support assumption.
Since $r\le\rho$ and $w\le1$, the inequality
\[
 \frac{(1+(321/500)^2)^2}{1-(321/500)^2}>(46/25)^2
\]
gives $|\beta|<321/500$.
Put $C=C_4$, $z=H/2$, $s=\sqrt{4/3}$.  The contact expression,
using the quantitative bound on the three-summand predecessor, obeys
\begin{align*}
 I_{\rm new}(y)\le{}&\mathcal H_4(z,y)+4\rho(K_3-C)
                    +4\rho(s-1)K_3\\
 &+C\{-|y|^3+\tfrac32\rho y^2+3\beta y-\rho/2\},
\mathcal H_4(z,y)={}&8\!\left\{\Phi\left(\frac{z-y/2}{\sqrt{3/4}}\right)-\Phi(z)\right\}
                 \\
 &+4y\phi(z)+z\phi(z)(y^2-1).
\end{align*}
Taylor expansion in the translation and the variance gives, uniformly
in $z,y$,
\[
 |\mathcal H_4(z,y)|<\frac{\phi(0)}6|y|^3+\frac{|y|}{5}+\frac1{15}.
\]
Here $\sup|\phi''|=\sup|(x\phi(x))'|=\phi(0)<2/5$ and
$\sup|x(3-x^2)\phi(x)|<3/5$ suffice.  For the last bound, on
$t=x^2\le3$ compare $e^t$ with its cubic Taylor polynomial and use
\[
 \tfrac94(1+t+t^2/2+t^3/6)-t(3-t)^2
 =\tfrac34\{7(t-9/14)^2+3/28\}+\tfrac58t^2(3-t)>0.
\]
For $t\ge3$, use $|x|^3\phi(x)<1/2$.
The second derivative of
$\Phi((z-y/2)/\sqrt{1-\lambda})$ has absolute value below $4/15$
on $0\le\lambda\le1/4$, giving the remainder $1/15$ after
multiplication by eight.

Discard $4\rho(K_3-C)<0$ and use $4(s-1)<5/8$.
For $R=|y|\ge4$, the remaining coefficient of $C$ is at most
$B(R)=-R^3+69R^2/25+963R/500+23/100$.
All coefficients of $B(4+t)$ are negative, so substituting $C>c_0$
is now valid.  Consequently
\[
 I_{\rm new}(y)<-\frac{643}{1875}R^3+
 \frac{17664}{15625}R^2+\frac{77257}{78125}R+\frac{7541}{46875}.
\]
The coefficients after translation by $R=101/25+t$ are all negative.
This excludes $|y|\ge101/25$ from the closed support by contact.
Finally $r^2\ge1/w$ gives $w>625/2116$.
\end{proof}



\subsection{The physical comparison at four summands}\label{j:finite:n4physical}
Write $F_{4,P}(H)=\Pp_P(S_4\le H)$, $T=4097/10000$, and
$R_*=17/4$. The preceding localization reduces a contrary extremizer
to $|X|\le R_*$, $\rho<46/25$, $625/2116\le w\le1$,
$|\beta|\le13/20$. The reference $Q$ and excess $\eta=\rho-r$
retain the original marginal; they do not replace it by a three-point law.

\paragraph{One measure carries all frequencies.}
For $\eta>0$ set $\nu=W\,dP/\eta$, tagged by
$(\sigma,y)=(\operatorname{sgn}X,|X|)$, and write $c_+=b,c_-=a$,
$p_0=1-w$, $g(q)=\varphi_Q(q)$. For
$r_q(\sigma,y)=\exp(i\sigma q\,\cdot)[0,c_\sigma,c_\sigma,y]$,
Hermite interpolation at $0,c,c$ gives
\begin{equation}\label{j:finite:n4measure}
 \varphi_P(q)-g(q)=\eta\int r_q\,d\nu,\qquad |r_q|\le |q|^3/6.
\end{equation}
The linear and quadratic terms agree on each sign class; the constant
terms agree after both classes and the zero atom are included.
Define $k_0=(c_\sigma^2y)^{-1}$, $k_2=(y-c_\sigma)^{-2}$ and
$k_1=(y-c_\sigma)^{-1}$ off the contacts, which have zero $\nu$ mass.
Multiplication by $W=y(y-c)^2$ and integration give
\[
 \eta\int k_0\,d\nu=p_0-P\{0\}\le p_0,\quad
 \eta\int_\sigma k_2\,d\nu=u-c_\sigma P\{\sigma c_\sigma\}\le u,\quad
 \int_\sigma k_1\,d\nu=0.
\]
Cauchy and the $k_2$ bound prove absolute integrability of $k_1$.
At $\eta=0$, $P=Q$ and no measure division is used. Also $a,b\le R_*$,
since $\E[y^2;\sigma]\le R_*\E[y;\sigma]$.
Thus for fixed nonnegative $\lambda_0,\lambda_\pm$ and signed
$\theta_\pm$, any whole-domain bound
\[
 S_H\ge\sup_{\sigma,\,0<y\le R_*}
 \{A_H-\lambda_0k_0-\lambda_\sigma k_2-\theta_\sigma k_1\}
\]
with $A_H=\Re\sum_j4\kappa_j(H)g(q_j)^3r_{q_j}$ gives
\begin{equation}\label{j:finite:n4priced}
 \eta\int A_H\,d\nu\le
 \Gamma_H(\eta):=p_0\lambda_0+u(\lambda_++\lambda_-)+S_H\eta.
\end{equation}
The signed frequency terms are summed before taking this supremum.
Neither $S_H$ nor $\Gamma_H$ is required to be nonnegative.

\paragraph{The nonlinear response.}
For $\psi=g+\delta$, the exact identity is
\[
 \psi^4-g^4=4g^3\delta+\delta^2(\psi^2+2\psi g+3g^2).
\]
If $|g|\le A\le1$, $|\psi|\le1$, $|\delta|\le E$, valid
remainder bounds are
\[
 E^2(p^2+2pA+3A^2),\quad p=\min(1,A+E),\qquad
 v(4A^2+2-v),\quad v=\min(E^2,2A^2+1).
\]
The second follows from
$12\delta^2\int_0^1(1-t)(g+t\delta)^2dt$ and
$|g+t\delta|^2=(1-t)|g|^2+t|\psi|^2-t(1-t)|\delta|^2$.
The sharper disk bound used by some cells is also finite algebra.
After rotating $g=a\ge0$, put $x=\Re\delta$, $s=|\delta|^2$;
the squared remainder equals
\[
 s^2\{36a^4+4a^2s+s^2+8ax(6a^2+s)+24a^2x^2\}.
\]
Convexity in $x$ leaves collinear endpoints or $|\psi|=1$.
With $r_+=\min(E,1)$, $m=\min(A,1-r_+)$, $r_-=\min(E,1+A)$,
the collinear modulus bounds are
\[
 r_+^2(6m^2+4mr_++r_+^2),\quad r_-^2(6A^2-4Ar_-+r_-^2),\quad
 \begin{cases}r_-^4,&r_-\le1,\\r_-^2(3r_-^2-8r_-+6),&r_-\ge1.
 \end{cases}
\]
On the unit circle maximize
$Q_A(s)=s^2[18A^4+12A^2+6-(12A^2+8)s+3s^2]$
on $[(1-A)^2,\min(E^2,(1+A)^2)]$: only its endpoints and
$s_-=(3A^2+2-A\sqrt{4-3A^2})/2$, if inside, are needed.
Convexity in $a^2$ reduces the remaining disk radii to $a=A$ or a
collinear endpoint. The largest collinear value and square root of
this circle maximum therefore bound the modulus; the minimum with
either preceding complete bound is valid.
On an excess cell $[L,U]$, an affine radius $e_j(\eta)\ge0$ is
used only when proved throughout the cell. Freezing
$p_j=\min(1,a_j+\max(e_j(L),e_j(U)))$ gives a quadratic remainder
with nonnegative leading coefficient. Freezing the entire radius
instead gives a valid constant disk bound. Bounds are combined as
complete alternatives, never coefficientwise.

\paragraph{From a finite frequency sum to every threshold.}
The closed-CDF majorant of Lemma~\ref{j:fourier:closed-majorant}, with
$\widehat K(t)=(1-t)+iR(t)$ and
$R(t)=(1-t)\cot(\pi t)+1/\pi$, gives
\[
 F_{4,P}(H)\le J_P(H/2),\quad
 J_P(x)=\tfrac12+\int_0^1
 \Re\{\widehat K(t)e^{-6itx}\varphi_P(3t)^4\}\,dt.
\]
This is a real compensated integral, including the closed boundary.
Put $Z=S_4/2$, $D_x=x-Z$, $S(t)=tR(t)$. Its integrand is
$(1-t)\E\cos(6tD_x)+S(t)\E[6D_x\operatorname{sinc}(6tD_x)]$.
The cotangent partial fractions give $0\le S\le1/3$,
$|S'|<2$, $|S''|<5$, and therefore
\[
 |F_x''(t)|\le42\E|D_x|+108\E D_x^2+24\E|D_x|^3.
\]
For the derivative constants, write
$S=\pi^{-1}-2\pi^{-1}(A+MG)$ with
$A=t^2/(1+t)$, $M=t^2(1-t)$ and
$G=\sum_{j\ge2}(j^2-t^2)^{-1}$.
On $[0,1]$, $(|A'|,|A''|,|M|,|M'|,|M''|)\le(3/4,2,1/6,1,4)$
and $(G,|G'|,|G''|)\le(3/4,10/27,106/135)$, by telescoping and
$j^2-t^2\ge3j^2/4$. For $t>1/2$, put $v=1-t$ and use $v\cot(\pi v)\le1/\pi$
to obtain $R(t)\ge0$; below $1/2$ positivity is immediate.
These estimates imply the stated bounds; the sinc
integral representation justifies differentiation by the third moment.

From \eqref{j:finite:n4-physical-moments} choose whole-cell bounds
\[
 \kappa_*\ge\min(\rho,|b-a|+\eta),\quad
 M_{4,*}\ge\min(R_*\rho,a^2-ab+b^2+(R_*+2\max(a,b))\eta).
\]
Then $D_2=x^2+1$, $D_4=x^4+6x^2+2|x|\kappa_*+(M_{4,*}+9)/4$
bound $\E D_x^2,\E D_x^4$, and Cauchy bounds the odd absolute moments.
For $N=256$, $t_j=(j+1/2)/N$, $q_j=3t_j$,
$\kappa_j(H)=\widehat K(t_j)e^{-iq_jH}/N$, set
\begin{gather*}
 C_H(Q)=1+2\Re\sum_j\kappa_j(H)g(q_j)^4-2\Phi(H/2),\\
 Q_H=\min\left\{\frac{687}{8N^2},
 \frac{42\sqrt{D_2}+108D_2+24\sqrt{D_2D_4}}{12N^2}\right\}.
\end{gather*}
This midpoint reserve bounds the original physical integral once.
The constant alternative uses $|x|\le2$, $D_2\le5$, $D_4<54$.
To bridge adjacent raw thresholds $H_\ell=-4+\ell/8$,
$0\le\ell\le57$, use
$P_j=\min(1,a_j+e_j+3/(2N))$, where $a_j\ge|g(q_j)|$ and
$e_j\ge|\varphi_P(q_j)-g(q_j)|$ on the excess cell.
The increment follows from $\E|X|\le1$.
If $d_j$ is the smaller of the exact integrals over the frequency cell of
$t^2(1-t)+t/3$ and $t^2(1-t)+t(1-t)/2$, then
$|J_P''|\le B=36\sum_jd_jP_j^4$.
Indeed both polynomials bound $t^2|\widehat K(t)|$; for the second,
retaining the first cotangent pole gives $S(t)<(1-t)/2$.
The single threshold-interpolation reserve is
$E_H=(B+1/4)/(16\cdot8^2)$, since $\sup|\Phi''|<1/4$.
The radii $\eta q_j^3/6$, $\max(a,b)^2p_0q_j^2$, and $1+a_j$
are all valid; the second uses
$y f[0,c,c,y]=f[c,c,y]-f[0,c,c]$, so $|r_q|\le q^2/y$, and
then the $k_0$ budget. Thus no contact division occurs.

\paragraph{The finite inequalities.}
Let $M_H$ bound $C_H(Q)-Tr$ on a geometry rectangle, and let
$R_2(\eta)$ bound the weighted nonlinear response on $[L,U]$.
The assembled bound is
\begin{equation}\label{j:finite:n4low}
 2\{F_{4,P}(H)-\Phi(H/2)\}-T\rho
 \le P_H(\eta):=M_H+Q_H+E_H+2\Gamma_H(\eta)+2R_2(\eta)-T\eta.
\end{equation}
It is convex quadratic in $\eta$, so its two nonpositive endpoints
suffice. The optional score bound checks both denominator faces
$46/25$ and $r_L+\eta$. For the latter the derivative numerator has
derivative $P_H''(\eta)(r_L+\eta)\ge0$, hence no interior maximum.
All functions defining $M_H,S_H,R_2$ are explicit finite expressions.
Their whole-cell enclosures have the following licenses.
For the reference, Taylor's theorem uses
$\partial_i C_H=8\Re\sum_j\kappa_jg^3\partial_i g$ and
$|\partial_{ik}C_H|\le8\sum_j|\kappa_j|(3D_iD_k+D_{ik})$.
A supporting plane bounds $r$ below: for
$f(\beta)=(1+\beta^2)/\sqrt{1-\beta^2}$,
$3ff''-(f')^2=(9+9\beta^2+15\beta^4-\beta^6)/(1-\beta^2)^3>0$
proves convexity of $r=w^{-1/2}f(\beta)$.
The response is evaluated without singular divided quotients by
\[
 r_q=\sum_{m=0}^{60}(i\sigma)^{m+3}b_mh_m+E_q,\quad
 |E_q|\le(q^3/6)(qR_*)^{61}/61!,
\]
where $b_0=q^3/6$, $b_m=b_{m-1}qR_* /(m+3)$,
$h_0=1$, $h_m=(y/R_*)h_{m-1}+(m+1)(c/R_*)^m$.
Hermite--Genocchi proves this expansion and
$|\partial_y r_q|\le q^4/24$, $|\partial_c r_q|\le q^4/12$.
The 273 nearest-node cells $y_k=k/64$, clipped to $[0,R_*]$,
therefore need radius $q^4/(48\cdot64)$.
The geometry correction uses
$|g^3-g_0^3|\le|g-g_0|(|g|^2+|g||g_0|+|g_0|^2)$.
The nonnegative prices are rational data. Each fixed choice is bounded
on every tagged cell; $k_0\ge1/(c_U^2y_U)$, while the other cost is
$-\lambda_\sigma z^2-\theta_\sigma z$, $z=1/(y-c)$.
Its endpoints and admissible vertex, on both reciprocal rays if the
cell meets a contact, bound it. Require $\theta_\sigma=0$ when
$\lambda_\sigma=0$. These rules prove the priced supremum over the
whole tagged domain. Factoring the difference of fourth powers also gives
$|\psi^4-g^4|\le E(p^3+p^2A+pA^2+A^3)$, $p=\min(1,A+E)$.
This absolute bound, or the full linear-plus-remainder expression, may be
used as a complete alternative.

For large excess, condition on $j$ nonnegative summands, assigning
zero to that class. With $p=P(X\ge0)$,
\[
 \mu_j=\frac{u(j-4p)}{2p(1-p)},\quad
 V_j=\frac14\left[j\left(\frac{v_+}{p}-\frac{u^2}{p^2}\right)
 +(4-j)\left(\frac{v_-}{1-p}-\frac{u^2}{(1-p)^2}\right)\right].
\]
Cantelli bounds each conditional CDF by $1$ if $x\ge\mu_j$ and by
$V_j/(V_j+(\mu_j-x)^2)$ otherwise. The closed sign-mass interval is
$w(1-\beta)/2\le p\le1-w(1+\beta)/2$.
On each threshold cell use the upper endpoint in Cantelli and the
lower endpoint in $\Phi$; degree-four Bernstein bounds retain the
joint binomial mixture. A sign-mass cell supplies numerator $d_i\ge0$
and correlated lower $r_i\le r$, giving score $d_i/(r_i+\eta)$.
With $r_B\le r$ on the parent, set
\begin{gather*}
 c_C^0=\max\{0,\max_i(d_i/T-r_i),2/(5T)-r_B\},\qquad
 c_C=10^{-4}\lceil10^4c_C^0\rceil+10^{-4},\\
 c_S=\max(0,979013/533000-r_B),\qquad c_B=\min(c_C,c_S).
\end{gather*}
Cantelli gives the first cutoff; the affine bound gives the second.
The two threshold rays $H\le-4$ and $H\ge25/8$ have score at most
$2/5$ by Cantelli and $\Phi(25/16)>4/5$. Thus $[c_B,\infty)$ is
covered, and it meets the low-excess interval at its included endpoint.

\begin{proposition}[Four-summand numerical assertion]\label{j:finite:n4numerical}
Divide $[625/2116,1]\times[0,13/20]$ into $20\times13$ closed
rectangles. A rectangle with $\min r^2\ge(46/25)^2$ is excluded.
Every other parent, or all four of its coordinate-bisected children,
passes the preceding large-excess bound and, on sixteen equal closed
subintervals of $[0,c_B]$, the low-excess inequalities at all 58
thresholds, their joins, both tails and both reflected orientations.
All 260 dispositions and all five parent/child request slots per parent
are verified by exact acceptance. Consequently $C_4\le T<\CE$.
\end{proposition}
\begin{proof}[Derivation from the numerical assertion]
Every localized contrary law belongs to a closed parent. Its reference
and common defect measure satisfy the derived inequalities on that parent
or one of its four children. The midpoint and threshold reserves cover
all central thresholds, and the two tails cover the rest.
Every $\eta\ge0$ belongs to the low or large-excess range, including
$\eta=0$ and the common endpoint. Hence its score is at most $T$,
contradicting $C_4>T$. Reflection accounts for the other discrepancy sign.
\end{proof}

\subsection{One support criterion for the remaining finite indices}
\label{j:finite:support}

The support calculations throughout the finite range use one elementary
criterion. Its inputs are a lower bound for the attained score, upper
bounds for moments of the same marginal, and a predecessor allowance.
The Gaussian calculation keeps the actual integer $n$ throughout.

\begin{lemma}[Direct Gaussian support criterion]
\label{j:finite:direct-gaussian-support-criterion}
Let $n\ge2$, and let $(P,H)$ be an attained positive closed-event
ratio maximizer with score $C$. Write
$\rho=\E|X|^3$, $\beta=\E[X|X|]$ and
$A=\sqrt{n-1}\Delta_{n-1}(P)/\rho$.
Suppose known numbers $c,r,b_-,b_+,S,e,q$ satisfy
\[
 0<c\le C,\quad 1\le\rho\le r,\quad
 -1<b_-\le\beta\le b_+<1,\quad q\ge0,
\]
\[
 \sqrt{n/(n-1)}\le S,\qquad
 n(\sqrt{n/(n-1)}-1)-1/2\le e,\qquad
 n(A-C)_+\le q.
\]
Fix $R\ge1$ and a support sign $\eta\in\{-1,1\}$, and put
$B_+=b_+$, $B_-=-b_-$. Define
\[
 Q_\eta(v)=-v^3+\tfrac32rv^2+3B_\eta v+re,
 \qquad
 L_\eta(v)=\tfrac25(v^2/2+1/2+e)+cQ_\eta'(v).
\]
Suppose that
\begin{equation}\label{j:finite:direct-support-ray-gates}
 Q_\eta(R),Q_\eta'(R),Q_\eta''(R),
 L_\eta(R),L_\eta'(R)<0,
 \qquad 1/5-3c<0.
\end{equation}
Define the actual-index Gaussian expression
\[
 G_n(H,t)=n\sqrt n\{\Phi((H-t)/\sqrt{n-1})-\Phi(H/\sqrt n)\}
 +nt\phi(H/\sqrt n)+\tfrac H2(t^2-1)\phi(H/\sqrt n),
\]
and suppose $U_\eta$ satisfies
\[
 \sup_{H\in\R}\eta G_n(H,R)\le U_\eta,
 \qquad U_\eta+cQ_\eta(R)+rSq<0.
\]
Then $\supp P$ contains no point in the ray
$\{\eta v:v\ge R\}$. The two signs can be checked separately,
with different signed-moment bounds.
\end{lemma}

\begin{proof}
Contact gives
\[
 0\le G_n(H,t)+
 C\{-|t|^3+\tfrac32\rho t^2+3\beta t+\rho e_n\}
 +\rho s_n n(A-C),
\]
where $s_n=\sqrt{n/(n-1)}$ and $e_n=n(s_n-1)-1/2$.
The last term is at most $rSq$. At $t=\eta v$ with $v\ge R$,
the bracket is at most $Q_\eta(v)$.
Since $Q_\eta'''=-6$, the first three gates make this polynomial
strictly negative on the full ray. Thus multiplication by $C\ge c$
gives an upper bound $cQ_\eta(v)$.

The derivative is
$n\phi(z)+Ht\phi(z)-ns_n\phi(s_ny)$, where
$z=H/\sqrt n$ and $y=(H-t)/\sqrt n$.
Taylor's theorem and $\|\phi''\|_\infty=\phi(0)$ bound
$|n\phi(y)-n\phi(z)-Ht\phi(z)|$ by $\phi(0)t^2/2$.
The identity $\partial_s[s\phi(sy)]=-\phi''(sy)$ bounds the
remaining scaling error by $n\phi(0)(s_n-1)$.
This proves, for every $n\ge2$,
$|\partial_tG_n(H,t)|\le\phi(0)(t^2/2+1/2+e)$.
Therefore the derivative in $v$ of
$G_n(H,\eta v)+cQ_\eta(v)+rSq$ is at most $L_\eta(v)$.
Its second derivative as a polynomial bound is
$L_\eta''=2/5-6c<0$; the last two endpoint gates imply
$L_\eta(v)<0$ on $[R,\infty)$.
Finally $G_n(H,-R)=-G_n(-H,R)$, so the assumed $U_\eta$
bounds $G_n(H,\eta R)$ uniformly in $H$.
The strict endpoint inequality and monotonicity contradict contact
on every point of the specified support ray.
\end{proof}

\paragraph{A finite certificate for the whole threshold line.}
Let $z=H/\sqrt n$ and $a=(H-R)/\sqrt{n-1}$. Direct
differentiation gives
\[
 \partial_HG_n(H,R)=ns_n\phi(a)+
 \left[-n-RH+\tfrac12(R^2-1)(1-H^2/n)\right]\phi(z).
\]
Choose a positive rational $H_0$ such that
$H_0\ge5\sqrt n$ and $H_0-R\ge5\sqrt{n-1}$.
An outward enclosure of the square roots makes this choice exact.
On the whole real line,
\[
 |\partial_H^2G_n(H,R)|\le M_{n,R}
 :=\frac{\sqrt n}{4}\left(\frac n{n-1}+1\right)
   +\frac25R+\frac{2|R^2-1|}{5\sqrt n}.
\]
On a closed threshold cell with midpoint $h$ and half-width $w$,
let $[g_-,g_+]$ and $[d_-,d_+]$ enclose $G_n(h,R)$ and
$\partial_HG_n(h,R)$. Taylor's theorem then gives the two upper
bounds
\[
 g_++\max(|d_-|,|d_+|)w+M_{n,R}w^2/2,
 \qquad
 -g_-+\max(|d_-|,|d_+|)w+M_{n,R}w^2/2.
\]
They bound $G_n(H,R)$ and $-G_n(H,R)$ respectively throughout
the cell. The certificate includes every consecutive cell from
$-H_0$ to $H_0$, including their endpoints.

On the two exterior rays, both distribution-function arguments lie
in the same normal tail, so
\[
 |G_n(H,R)|\le n\sqrt n\,\tau+
 \left(nR+\tfrac12H_0|R^2-1|\right)\phi(H_0/\sqrt n),
 \quad |H|\ge H_0,
\]
where $\tau=3\cdot10^{-7}>\Phi(-5)$ is sufficient.
Take the largest positive-cell bound and the exterior bound for
$U_+$, and the largest negative-cell bound and the exterior bound
for $U_-$. This proves a statement about the entire threshold line,
not just about its grid points.



\begin{proposition}[Actual-index support rows]\label{j:finite:supportrows}
Under $C_{n-1}\le\CE<C_n$, each row below bounds the support of the
same positive closed maximizer by $(-R,R)$. Repeated index ranges mean
one check at every actual integer, not a transfer between integers.
\begin{center}\small
\begin{tabular}{c|c|c}
$n$&third-moment interval&$R$\\\hline
5&$[1,11/10],[11/10,23/20],[23/20,237/200],[237/200,6/5]$
 &$11/5,47/20,49/20,99/40$\\
5&$[6/5,5/4],[5/4,7/5],[7/5,3/2]$&$13/5,117/40,25/8$\\
5&$[3/2,8/5],[8/5,17/10],[17/10,7361/4000]$&$133/40,7/2,15/4$\\
6--14&$[1,11/10],[11/10,23/20]$&$23/10,12/5$\\
6--9&$[23/20,237/200],[23/20,6/5]$&$5/2,13/5$\\
6--9&$[6/5,7/5],[7/5,3/2],[3/2,8/5]$&$3,16/5,17/5$\\
6--9&$[8/5,17/10],[17/10,7361/4000]$&$18/5,19/5$\\
10--14&$[23/20,71/50]$&$3$
\end{tabular}
\end{center}
The last five lines comprise all 26 low-moment and 25 complementary
actual-index rows. Their endpoint tests are instances of the same
criterion, including both support rays and all real thresholds.
\end{proposition}
\begin{proof}[Mathematical meaning and finite inequalities]
Use $c=819/2000<\CE$, $q=0$, and $e\le91/1000$.
For each upper moment endpoint $r$, the signed-moment bound supplies
$|\beta|\le b$, where $b$ is the outward rational root of
$(1+b^2)^2\ge r^2(1-b^2)$ recorded in the support certificate.
The certificate verifies the five strict polynomial gates and the strict
Gaussian endpoint test. For $n=5$ use $H_0=16$; for $6\le n\le9$
use $H_0=35/2$, enlarged to 18 for the complementary rows at $n=9$;
for $10\le n\le14$ use $H_0=22$. All threshold meshes have width
$1/10$. The value at each endpoint, enlarged by $M_{n,R}/800$,
bounds both signs on the intervening cell: the interpolation error for
$|f''|\le M$ is $M h^2/8$, by comparison with
$\pm M(x-a)(b-x)/2$. The displayed exterior-ray argument applies
because $H_0^2\ge25n$ and $(H_0-R)^2\ge25(n-1)$.
All numerical comparisons are strict rational inequalities with the
enclosures of Section~\ref{j:cert:arithmetic}. Thus the preceding support
criterion applies before any of these radii enters a Fourier estimate.
Additional actual-index rows at $15$--$39$, and the separate two-endpoint
support test at $40$--$42$, use the same contact license; their complete
input domains are part of the numerical propositions below.
\end{proof}

\subsection{Finite-frequency refinements}
\label{j:finite:small-fourier}
At a small index, replacing a finite frequency integral by an integral
over the whole half-line can lose much of the available margin. The
following estimates retain the finite interval and the actual integer
power. They use the moments of the same supported marginal throughout.

\subsubsection{A finite-interval compensated estimate}
Let $n>2$ be an integer, $\E X=0$, $\E X^2=1$, $|X|\le R\le6$,
$\gamma=\E X^3$, $\mu_j=\E X^j$, and $f(t)=\E e^{itX}$.
Assume $\mu_4\le M\le111/10$ and $|\gamma|\le g\le37/20$.
The simultaneous coefficient and moment bounds in the signed moment atlas
of Section~\ref{j:local:moments} give
\[
 |a|\le A_*,\quad |z|\le Z,\quad \mu_6\le M_6,
 \qquad a=(3-\mu_4)/24,\quad z=(10\gamma-\mu_5)/120.
\]
Define
\[
 E=\frac{M_6}{720}+\frac{g^2}{72}
       +\frac{(1/2+g^2/144)^3}{5},\qquad
 B=A_*+Z/2+E/4,
\]
\[
 a_n=(n-1)/2,\qquad b_n=(n-2)(45-M)/96,\qquad
 I_j(A)=\int_0^{1/2}t^j e^{-At^2}\,dt.
\]
All these quantities are nonnegative, and the two damping rates are
positive. Put
\begin{align*}
 L_0={}&\sqrt n\left\{nZI_2(a_n)
       +\frac{n(n-1)A_*g}{6}I_4(a_n)
       +\frac{n^3g^3}{1296}I_6(n/2)\right\},\\
 L_H={}&\sqrt n\left\{nA_*I_2(a_n)
                              +\frac{n^2g^2}{72}I_4(n/2)\right\},\\
 L_2={}&n\left\{nEI_3(a_n)+\frac{n(n-1)B^2}{2}I_5(b_n)\right\}.
\end{align*}
For every real raw threshold $H$, the complete low-frequency estimate is
\begin{equation}\label{j:finite:finite-low}
 \int_0^{1/2}\frac{\left|\Im\left[e^{-itH}
 \{f(t)^n-e^{-nt^2/2}(1-in\gamma t^3/6)\}\right]\right|}{t^3}\,dt
 \le\frac{L_0+L_H|H|}{\sqrt n}+\frac{L_2}{n}.
\end{equation}

To prove this finite-interval version, use the Taylor decomposition
\eqref{j:fourier:taylor} with $h(t)=e^{-t^2/2-i\gamma t^3/6}$.
The same polynomial-power remainder gives
$|f^n-h^n-n(f-h)h^{n-1}|
\le n(n-1)B^2t^8e^{-b_nt^2}/2$.
Retain the integration interval $[0,1/2]$ in that proof.
The leading rotated imaginary terms are bounded by
$ne^{-a_nt^2}[(Z+A_*|H|)t^5+(n-1)A_*gt^7/6]$;
the cubic Gaussian phase remainder contributes
$e^{-nt^2/2}(|n\gamma t^3/6|^3/6+
|H|t(n\gamma t^3/6)^2/2)$.
Division by $t^3$ and integration give exactly $L_0,L_H,L_2$ above.
The damping and all Taylor hypotheses are those of the same proof;
no half-line extension or asymptotic replacement is made.

For the small indices the quantities $I_j(A)$ are enclosed directly.
The degree-20 and degree-21 Taylor polynomials of $e^{-x}$ are,
respectively, upper and lower bounds on all $x\ge0$, by the sign of
the integral remainder. Their integrated values are the rational sums
\[
 \sum_{k=0}^{m}\frac{(-A)^k}{k!(j+2k+1)2^{j+2k+1}},
 \qquad m=20,21.
\]
The selected routines require $0\le A\le8$ and $0\le j\le16$;
the rates above meet those gates for $5\le n\le9$.
Every coefficient multiplying an integral is nonnegative. When two
complete triples $(L_0,L_H,L_2)$ are available, either may be selected
by a rule independent of $H$. Selecting different coefficients from
different triples is not justified by this argument.

\subsubsection{Containing rectangles and exact finite powers}
The supported fourth-moment bound also improves the real coordinate.
If $\rho\in[\rho_l,\rho_u]$, $|X|\le R$, and $\mu_4\le M_4$, then
\[
 \Re f(t)\le1-t^2/2+M_4t^4/24.
\]
If $R^2t^2\le30$, a second lower endpoint is
\[
 \Re f(t)\ge1-t^2/2
              +\rho_l^2t^4(1-R^2t^2/30)/24.
\]
These follow from the global fourth-degree upper and sixth-degree
lower Taylor polynomials of cosine, $\mu_6\le R^2\mu_4$, and
$\mu_4\ge\rho^2$ by Cauchy's inequality. The nonnegative coefficient
gate is needed before substituting $\rho_l^2$. Intersect these real
intervals with the anchored bounds of Lemma~\ref{j:fourier:anchor}.

\begin{lemma}[Two corners of a complex-power rectangle]
\label{j:finite:two-corners}
Let $n\ge2$ be an integer, $a\ge0$, $0<x_0\le x_1$ and $Y\ge0$.
If
\[
 (n-2)Y\le3x_0/2,\qquad (2n-1)x_0^n\ge(n-1)a,
\]
then for $x\in[x_0,x_1]$ and $|y|\le Y$,
\[
 |(x+iy)^n-a|^2
 \le\max\{|(x_0+iY)^n-a|^2,|(x_1+iY)^n-a|^2\}.
\]
\end{lemma}
\begin{proof}
Write $z=x+iy$, $r=|z|$ and $F(x,y)=|z^n-a|^2$.
Since $|\arg z|\le Y/x_0$, the first gate and $3/2<\pi/2$ imply
$\Re z^{n-2}\ge0$. Direct differentiation gives
\begin{align*}
 F_{yy}&=2nr^{2n-4}\{r^2+2(n-1)y^2\}
                         +2an(n-1)\Re z^{n-2}\ge0,\\
 F_{xx}&=2nr^{2n-4}\{r^2+2(n-1)x^2\}
                         -2an(n-1)\Re z^{n-2}.
\end{align*}
Divide the latter by $2nr^{n-2}>0$. Its first term is
$r^{n-2}((2n-1)x^2+y^2)\ge(2n-1)x_0^n$; the subtracted
term is at most $a(n-1)$. Thus $F$ is convex in each coordinate.
Apply convexity first in $y$ and then in $x$, and use
$F(x,Y)=F(x,-Y)$.
\end{proof}

For Fourier smoothing take $a=e^{-nt^2/2}$ and any proved containing
rectangle for $f(t)$. Both gates must hold throughout the frequency
cell. If either fails, the geometric-sum estimate remains available.
The maximum of the two pointwise corner bounds is enclosed before
integration; the maximum of their separate integrals would not suffice.
Complex powers are finite polynomial recurrences, so zero or equal
bases create no division. A fourth-Taylor endpoint is kept fixed on
each integration cell when differentiating its smooth majorant.

Cancellation can be retained when computing the real endpoints.
For example, on $0\le u\le1$,
\[
 1-u^2/2+M_4u^4/24-e^{-u^2/2}
 =\frac{M_4-3}{24}u^4
   +\sum_{k=3}^{10}\frac{(-1)^{k+1}u^{2k}}{2^kk!}+R(u).
\]
For derivative order $j=0,1,2$ with $u=ct$, the absolute remainder
is at most $2(22)_j u_{\max}^{22-j}c^j/(2^{11}11!)$; consecutive
absolute tail terms have ratio below $1/2$ on this domain.
Alternatively, at higher frequencies enclose each explicit smooth
residual $d(u)=\cos u-e^{-u^2/2}$ or
$d(u)+(\rho_u-1)u^3/6$ by its midpoint value plus the half-width
times an enclosure of its derivative on the entire cell.
Intersecting two such valid enclosures preserves their common function.
These operations differentiate scalar envelopes, not the law's CDF.


\subsubsection{A sine improvement on the acquired $26$--$30$ region}
\label{j:finite:finite-sine}
Let the positive and negative support magnitudes lie in
$[a_+,R_+]$ and $[a_-,R_-]$, respectively. Assume
\[
 0<A\le a_+^2+a_-^2,\qquad
 B\ge R_+^2+R_-^2,\qquad M\ge\max(R_+,R_-).
\]
Put $R(u)=1-Au^2(1-Bu^2/42)/20$.
For $u\ge0$, $Mu\le\pi$ and $Bu^2\le42$,
\begin{equation}\label{j:finite:finite-sine-bound}
 |\Im f(u)|^2\le(\rho^2-\zeta^2)u^6R(u)^2/36.
\end{equation}
In particular $R(u)\ge1-21A/(40B)\ge19/40$.
Outside these frequency gates the ordinary sine-circle inequality
remains available.

Here is a proof retaining the same marginal. Write
$\operatorname{sinc}t=\int_0^1\cos(vt)\,dv$. For $u\sqrt z\le\pi$,
\[
 0\le-\frac{d}{dz}\operatorname{sinc}(u\sqrt z)
 =\frac{u^2}{2}\int_0^1v^2\operatorname{sinc}(uv\sqrt z)\,dv
 \le\frac{u^2}6-\frac{u^4z}{60}+\frac{u^6z^2}{1680}.
\]
The upper bound follows by integrating the global polynomial
$\operatorname{sinc}t\le1-t^2/6+t^4/120$; nonnegativity follows
from $0\le uv\sqrt z\le\pi$.
Integrating between $a^2$ and $b^2$ gives
\[
 |\operatorname{sinc}(au)-\operatorname{sinc}(bu)|
 \le\frac{|a^2-b^2|u^2}{6}
 \left(1-\frac{(a^2+b^2)u^2}{20}
       +\frac{(a^4+a^2b^2+b^4)u^4}{840}\right).
\]
The last bracket is at most $R(u)$, since
$a^4+a^2b^2+b^4\le B(a^2+b^2)$ and
$1-Bu^2/42\ge0$. At $a=b$ use the undivided inequality.

For the centered two-point law at $a,-b$, put
\[
 e_{a,b}=\frac{2a^2b^2}{a+b},\quad
 r_{a,b}=\frac{ab(a^2+b^2)}{a+b},\quad
 h_{a,b}=\frac{abu}{a+b}
       \{\operatorname{sinc}(au)-\operatorname{sinc}(bu)\}.
\]
The identity $(2ab)^2+(a^2-b^2)^2=(a^2+b^2)^2$ gives
$\|(e_{a,b},6h_{a,b}/(u^3R(u)))\|_2\le r_{a,b}$.
Integrate against the exact centered two-point mixture from the
ordinary sine-circle proof. The three integrals are $\zeta$, $\rho$
and $\Im f(u)$, respectively. The triangle inequality proves
\eqref{j:finite:finite-sine-bound}. The case $u=0$ is immediate.

For each of the $55$ acquired rows ($26\le n\le30$,
$32\le j\le42$), retain its closed probability interval
$[p_l,p_u]$, an actual-noise bound $s\le s_u$, and an upper
bound $d>0$ on the inward sum $a_0+a_1$.
The interface requires $d<p_l$, $d\le1-2p_u$ and
\begin{equation}\label{j:finite:b1-support-map}
 \frac{p_u^2+(1-p_u-d)^2}{p_u(1-p_u)+s_u}\ge\frac{17}{10},
 \qquad R_-^2+R_+^2\le\frac{1037}{100},
 \qquad \max(R_-,R_+)\le\frac{13}{5}.
\end{equation}
Indeed the minimum of $(p-a_0)^2+(1-p-a_1)^2$ under
$a_0,a_1\ge0$, $a_0+a_1\le d\le1-2p$ is
$p^2+(1-p-d)^2$: use the full inward sum and differentiate in
$a_0$, obtaining $2(1-2p-d)+4a_0\ge0$.
This numerator decreases in $p$ on the retained interval, while
$p(1-p)+s$ increases. Thus \eqref{j:finite:b1-support-map} gives the
required standardized inner support sum. It licenses
\eqref{j:finite:finite-sine-bound} with
$A=17/10$, $B=1037/100$, $M=13/5$ on every row.

\paragraph{A direct resource-to-sine map at indices $26$--$30$.}
The three common constants in \eqref{j:finite:b1-support-map} follow
directly from the acquired floor-$25$ resource.
Write its original simultaneous intercepts as $d_0,L$, and its prices
as $b,g,h$, so that
\[
 d_{\rm floor}-b=g d_0=hL>0,\qquad b\ge0.
\]
These identities are checked exactly against the exported residual;
$L$ is the original joint intercept, even when another bound further
clips the variance. The weighted resource implication gives
\[
 \frac n{25}(b+g\delta)+h\lambda_n\le d_{\rm floor},
 \qquad
 \frac n{25}\frac{\delta}{d_0}+\frac{\lambda_n}{L}\le1.
\]
The second inequality follows by subtracting $nb/25$, dividing by
$d_{\rm floor}-b$, and using $n\ge25$. Therefore
\[
 \delta\le \delta_u:=\min\left\{\frac{25d_0}{n},
                 \frac{25d_{\rm floor}/n-b}{g}\right\}.
\]
The bounded conditional supports give $s\le\mathcal B\delta$, where
$\mathcal B=\max(O_0,O_1)\max(p_u,1-p_l)$ and $\lambda_n=ns$.
Substitution in the joint plane yields
\[
 \lambda_n\le
 \left\{L^{-1}+(25\mathcal B d_0)^{-1}\right\}^{-1}=:\lambda_u.
\]
This step uses one joint plane before optimizing. For all $55$ rows,
the choices $d=\delta_u$ and $s_u=\lambda_u/n$ satisfy the two
preceding gates and
\[
 \frac{p_u^2+(1-p_u-\delta_u)^2}
      {p_u(1-p_u)+\lambda_u/n}>\frac{17}{10}.
\]
The smallest left side is greater than $1716/1000$ and occurs
at $n=26$ on $P_{38}$. The same rows satisfy
$R_-^2+R_+^2\le1037/100$ and $\max(R_-,R_+)\le13/5$.
The actual-index support proofs and the full $100$-panel resource
cover are checked separately. Thus the finite-sine majorant applies
on every surviving panel. Its integral uses only these common
constants and the rational moment and frequency intervals below.

Here $\eta,\zeta$ denote the symmetrization quantities of
Lemma~\ref{j:fourier:surrogate}, not the four-summand excess. Their
bounds follow directly from that lemma. Since
$(13/5)^2< (11/10)^2+8$, the general bounded-support inequality
for $\eta$ gives $\eta\le b(\rho)$ throughout $[11/10,6/5]$,
where $b(r)=(\sqrt{r^2+8}-r)/2$.
Lemma~\ref{j:fourier:surrogate} supplies $\zeta\ge b(\rho)$ and the supported bound for $\eta$.
Since $b$ decreases and $r+b(r)$ increases, on a strip $[r_l,r_u]$
we may use simultaneously
\[
 \rho+\eta\le r_u+b(r_u),\qquad
 \rho^2-\zeta^2\le r_u^2-b(r_u)^2.
\]
Rational brackets $z_l,z_u$ for $b(r_u)$ are checked by
$z_l^2+r_uz_l\le2\le z_u^2+r_uz_u$ with $0<z_l\le z_u$.
The integrator uses $r_u+z_u$ and $r_u^2-z_l^2$ in their indicated
upper-bound directions. The value $z_u$ is not separately asserted
to bound $\eta$ at the lower end of the strip.

The selected $23$ closed strips cover $[11/10,6/5]$. All five
actual indices are evaluated on each strip, giving $115$ complete
four-term Fourier inequalities. The optional fourth-Taylor real
endpoint and Lemma~\ref{j:finite:two-corners} are used only when their
whole-cell gates hold. The remaining cells use a previously valid
geometric-sum or direct interval bound.

The implementation sometimes integrates the maximum of the two
corner majorants more accurately. If $f_0,f_1$ are twice continuously
differentiable on $[t_0-h,t_0+h]$, enclose their midpoint values above
by $a_i$, their midpoint derivatives by $[d_i^-,d_i^+]$, and both
second derivatives above by $C\ge0$. Taylor's theorem bounds the
integral of their pointwise maximum by
\[
 \int_0^h\max_i(a_i+d_i^+x)\,dx
 +\int_0^h\max_i(a_i-d_i^-x)\,dx+\frac{Ch^3}{3}.
\]
Each affine maximum has at most one crossing, found exactly. This
bounds the maximum before integration and assumes nothing about
how often the original smooth curves cross.

\subsection{A common count comparison}\label{j:finite:finite-counts}
The finite argument follows the number of positive summands. Conditional on
that number, the sum is a binomial lattice point plus a centered residual.
The residual can move mass across a threshold in both directions. Keeping
its outgoing and incoming contributions together is essential near a
binomial jump. Everything in this section concerns one common marginal;
no finite-atom assumption is made.

\subsubsection{Exact conditioning and the common moments}
Use the sign normalization \eqref{j:core:normalization}, writing
$h=1/\ell=pq/m$ and $v=pq$. Thus $hX=B-p+U$,
$h^2=v+s$, $\lambda=ns$ and
$w_0=\lambda(1-\alpha)/(nq)$, $w_1=\lambda\alpha/(np)$.
At $s=0$ both conditional residuals vanish and any $\alpha\in[0,1]$
may be used; zero atoms belong to $B=0$ as before.

Write $W_j$ for a sum of $n-j$ independent copies of $U\mid B=0$
and $j$ independent copies of $U\mid B=1$, all independent. For
$b_j=\binom njp^jq^{n-j}$ and $z=np+hH=k+u$,
\begin{align}
 \Pp(S_n\le H)&=\sum_{j=0}^n b_j\Pp(W_j\le k+u-j),
 \label{j:finite:count-identity}\\
 \Pp(S_n\le H)&\le F_{\mathrm{Bin}(n,p)}(k)
 -b_k\Pp(W_k>u)+\sum_{j>k}b_j\Pp(W_j\le u-(j-k)).
 \label{j:finite:count-bound}
\end{align}
The second line discards nonnegative outgoing probabilities at $j<k$.
The outgoing event is strict and the incoming event is closed. In particular,
an atom exactly on the threshold is retained on the CDF side.
Also
\[
 V_j:=\E W_j^2=(n-j)w_0+jw_1
 =\lambda\left(1+\frac{(j-np)(\alpha-p)}{npq}\right).
\]

Suppose support acquisition gives
$U_0\in[-O_0,a_0]$ and $U_1\in[-a_1,O_1]$, with nonnegative
endpoints. For a centered $T\in[-L,R]$ of variance $w$,
\begin{equation}\label{j:finite:conditional-moments}
 w\le LR,\quad w^2/L-Lw\le\E T^3\le Rw-w^2/R,
 \quad \E T^4\le(R-L)\E T^3+LRw.
\end{equation}
The first inequality integrates $(T+L)(R-T)\ge0$. The two cubic
bounds integrate $(T+L)(T-w/L)^2\ge0$ and
$(R-T)(T+w/R)^2\ge0$. Multiplying the first polynomial by $T^2$
gives the fourth-moment bound. If $L=0$ or $R=0$, centering forces
$T=0$; the divisions are replaced by this deterministic case.

The common cubic identities \eqref{j:core:conditional-moments}, in
these coordinates, read $\rho h^3=v(1-2v)+3v(w_0+w_1)-q\mu_0+p\mu_1$,
$\gamma h^3=v(1-2p)+3v(w_1-w_0)+q\mu_0+p\mu_1$, and
$\beta h^2=v(1-2p)+s(2\alpha-1)$. The same pair
$(\mu_0,\mu_1)$ is retained in every probability bound below.

\subsubsection{Outgoing and incoming probability bounds}
A useful outgoing minorant is
\[
 G_a(x)=x(x+a)^2(2a-x),\qquad a>0.
\]
It is nonpositive outside $[0,2a]$, and its maximum is
$K a^4$, where $K=9/4+3\sqrt3/2$. Consequently
\[
 \Pp(W_k>u)\ge \frac{\E G_a(W_k-u)}{K a^4}.
\]
This is a global polynomial inequality, including $W_k=u$.
Its expectation uses
\[
 \E W_k^3=(n-k)\mu_0+k\mu_1,
\]
\[
 \E W_k^4=(n-k)\E U_0^4+k\E U_1^4
 +3\{V_k^2-(n-k)w_0^2-kw_1^2\}.
\]
Thus \eqref{j:finite:conditional-moments} makes the outgoing bound affine in
the two common cubic moments. Negative lower bounds remain valid; they
are not mistaken for negative probabilities.

For an incoming count $j>k$ and phase $u\le1/2$, put $d=j-k-u>0$.
Cantelli gives
\[
 \Pp(W_j\le-d)\le\frac{V_j}{V_j+d^2}.
\]
For completeness, for a centered $Y$ of variance $V$ and $a>0$,
$\Pp(Y\ge d)\le\E(Y+a)^2/(d+a)^2=(V+a^2)/(d+a)^2$.
Choosing $a=V/d$ proves the formula; $V=0$ is deterministic.
Exponential bounds can be sharper. If $\E Y=0$, $Y\le B$, $B>0$,
and $\E Y^2\le w$, then for $t\ge0$,
\begin{equation}\label{j:finite:one-sided-mgf}
 \E e^{tY}\le
 M(w;B,t):=\frac{B^2e^{-tw/B}+we^{tB}}{B^2+w}.
\end{equation}
Indeed, the quadratic interpolating $e^{tx}$ at $-w/B$ with multiplicity
two and at $B$ majorizes it on $x\le B$: the Hermite remainder has
sign $(x+w/B)^2(x-B)\le0$. Its quadratic coefficient is nonnegative.
The centered two-point law on $\{-w/B,B\}$ with variance $w$ evaluates
its expectation exactly. This proves the formula also with an upper
variance bound. At zero variance use $Y=0$.
Differentiation shows that $M$ is increasing and concave in $w$:
with $x=t(B^2+w)/B$,
\[
 M''(w)=\frac{B^2e^{-tw/B}}{(B^2+w)^3}
          (x^2+2x+2-2e^x)\le0.
\]
Also
\[
 M'(w)=\frac{B^2}{(B^2+w)^2}
 \{e^{tB}-e^{-tw/B}(1+x)\}\ge0,
\]
since $e^x\ge1+x$. Its tangents therefore give affine upper bounds throughout the variance
half-line. Applying \eqref{j:finite:one-sided-mgf} to $-U_0,-U_1$ and
multiplying the corresponding factors bounds $\E e^{-tW_j}$.
Chernoff then gives $e^{-td}\E e^{-tW_j}$ as an incoming upper bound.

One can also retain the common cubic moments. On $[A,B]$, $A<0<B$,
the function
\[
 r_t(x)=\frac{e^{tx}-1-tx}{x^2}
       =t^2\int_0^1(1-r)e^{trx}\,dr
\]
is increasing and convex. Its chord gives
$e^{tx}\le1+tx+c_2x^2+c_3x^3$, with $c_2,c_3\ge0$.
Applied to $-U_i$, choose $C_2\ge c_2$ and
$0\le C_{3,-}\le c_3\le C_{3,+}$, and put
\[
 Q_i(\mu_i)=\max\{0,1+C_2w_{i,+}-C_{3,-}\mu_i,
                     1+C_2w_{i,+}-C_{3,+}\mu_i\}.
\]
This is a nonnegative convex upper bound for its MGF. The maximum with zero
also makes it nonnegative at infeasible vertices of a relaxed moment polygon.
For fixed positive pivots $P_0,P_1$, weighted AM--GM yields
\[
 Q_0^{n-j}Q_1^j\le P_0^{n-j}P_1^j
 \left(\frac{n-j}{n}\frac{Q_0}{P_0}
       +\frac{j}{n}\frac{Q_1}{P_1}\right)^n.
\]
The right side is convex in $(\mu_0,\mu_1)$. A fixed exponential rate
and its phase chord produce a convex common-moment bound on an entire
closed phase interval. Choose a strategy on that whole interval and
moment polygon before comparing different strategies; a pointwise
minimum of convex functions need not be convex.

\subsubsection{Normalization and a finite closed cover}
Let $\mathcal M$ be a compact polygon containing the actual common
cubic moments, obtained from proved moment and defect restrictions.
Suppose $F$ is convex on $\mathcal M$ and $L$ is affine and strictly
positive there. Then
\begin{equation}\label{j:finite:normalized-vertex}
 \frac{F(x)}{L(x)}\le
 \max_{v\text{ vertex of }\mathcal M}\frac{F(v)}{L(v)}.
\end{equation}
To see this, write $x=\sum_v\theta_vv$. Convexity bounds the numerator
by $\sum_v\theta_vF(v)$; affinity identifies the denominator with
$\sum_v\theta_vL(v)$. The ratio is therefore bounded by the largest
vertex ratio. Segments and single points are included. If the lower
bound for $\rho$ is $\max\{r_0,L\}$, split $\mathcal M$ along
$L=r_0$ and retain both closed pieces.

\begin{proposition}[Whole-cell count certificate]\label{j:finite:count-cell}
Fix a closed box of acquired parameters, a count $k$, and a closed
phase interval. Let $B_k$ bound the binomial CDF above, $b_k^-$ the own
mass below, and $b_j^+$ the incoming masses above. Let $L_k$ be a
proved outgoing lower bound, $A_j\ge0$ proved incoming upper bounds,
and $N$ a lower bound for $\Phi((k+u-np)/(\sqrt n h))$.
Suppose the combination
\[
 B_k-b_k^-L_k+\sum_{j>k}b_j^+A_j-N
\]
is bounded on the whole phase interval by convex Bernstein control
functions $F_i(\mu_0,\mu_1)$. Let $r>0$ be the affine lower bound for
$\rho$ on each closed normalization piece and let $n_+\ge\sqrt n$.
Then the positive normalized CDF error is at most
\[
 n_+\max_{i,\,v\text{ vertex}}
       \frac{\max\{0,F_i(v)\}}{r(v)}.
\]
\end{proposition}
\begin{proof}
Insert the probability bounds into \eqref{j:finite:count-bound}. For the outgoing
term, $b_k\Pp(W_k>u)\ge b_k^-\Pp(W_k>u)\ge b_k^-L_k$,
even if $L_k<0$. The nonnegativity of the incoming bounds justifies the
upper mass substitutions. The Bernstein basis is nonnegative and sums
to one. Take the positive part and apply \eqref{j:finite:normalized-vertex};
the positive part of a convex function is convex. Finally enlarge
$\sqrt n$ to $n_+$.
\end{proof}

The acquired $p,h,H$ bounds must first put every actual threshold into
one of the original roots $k=0,\ldots,n$, $u\in[-1/2,1/2]$.
If they do not, prove the remaining phase tails separately. A tied
nearest count belongs to both closed roots. Apply
Proposition~\ref{j:cert:cover} only after every leaf is either a
proved necessary-condition contradiction or satisfies the whole-cell
score bound. The finite tree then covers infinitely many parameters.
Its leaves never assert that every feasible moment tuple is an iid law;
they need only contain all actual laws.


\subsubsection{Necessary filters for the same count domain}\label{j:finite:filters}
In the following formulas write $s_i=w_i$, $x=\mu_0$, $y=\mu_1$.
The preceding cubic identities and support-moment inequalities give all
three quantities for the same law; they may not be chosen separately.
These identities restrict the single pair $(x,y)$ to a rectangle and
linear strips. For example, replacing $p$ by the midpoint $p_m$ in
$-qx+py$ incurs error $(p-p_m)(x+y)$. On a closed cubic rectangle
its absolute value is at most the probability half-width times
$\max(|x_l+y_l|,|x_u+y_u|)$. In the signed-cubic identity the
corresponding factor is $y-x$, bounded by
$\max(|y_l-x_u|,|y_u-x_l|)$. All other products use directed interval
endpoints. Thus these strips contain every actual common pair.

\subsubsection{The physical defect as a further half-plane}
Let $\sigma=\sqrt{10}$ and retain the physical gap defect
$D=(3+\sigma)\rho-\gamma-3J$.
The endpoints of the central gap are
$A=(p-a_0)/h$, $B=(q-a_1)/h$.
The outer-gap inequality gives
\[
 D\ge\sigma\rho-\gamma-6mAB
       +6B(\E X_+^2-mB)+6A(\E X_-^2-mA).
\]
Substituting the count coordinates and the two cubic identities yields
\begin{align}
 Dh^3\ge{}&B_\sigma(p)+6vF(a_0,a_1)\notag\\
 &+qs_0\{3p(\sigma+3)-6a_0\}
  +ps_1\{3q(\sigma+1)-6a_1\}\notag\\
 &-(\sigma+1)qx+(\sigma-1)py,\label{j:finite:defect-halfplane}\\
 B_\sigma(p)&=v\{(\sigma-1)+2p-(2\sigma+6)v\},\notag\\
 F(a_0,a_1)&=a_0+a_1-a_0^2-a_1^2-a_0a_1.\notag
\end{align}
The polynomial $B_\sigma$ is the Bernoulli defect square.
Replacing $\sigma$ by a lower enclosure decreases the complete right
side, because its coefficient is exactly $\rho h^3\ge0$.
After that replacement, signed interval products and the explicit
nonnegative-coefficient gates license the termwise lower bounds.
If such a gate fails, this additional half-plane is omitted.

Substitution of \eqref{j:finite:conditional-moments} into
\eqref{j:finite:defect-halfplane} gives the stronger scalar inequality
\begin{align}
 Dh^3\ge{}&B_\sigma+6vF(a_0,a_1)\notag\\
 &+q\left\{[3p(\sigma+3)-(\sigma+7)a_0]s_0
                         +(\sigma+1)s_0^2/a_0\right\}\notag\\
 &+p\left\{[3q(\sigma+1)-(\sigma+5)a_1]s_1
                         +(\sigma-1)s_1^2/a_1\right\}.
 \label{j:finite:defect-sharp}
\end{align}
The terms with $a_i=0$ are treated by $s_i=0$, not by division.
Also $s_i\le O_i a_i$, by averaging the nonnegative quadratic
$(U_i+O_i)(a_i-U_i)$, with reflection for class one.

On a rectangle clipped by $a_0+a_1\le d<1$, the function $F$ is
nonnegative and concave: its Hessian is
$\left(\begin{smallmatrix}-2&-1\\-1&-2\end{smallmatrix}\right)$.
Its minimum therefore occurs at a vertex of the clipped polygon.
There is no requirement $d<1/2$ for this vertex argument.
The older coordinate-monotonicity shortcut requires separate gates.
For example, a lower bound for the $a_0$ derivative in
\eqref{j:finite:defect-sharp} is
\[
 6v_l(1-2d)-q_u(\sigma_u+7)(s_0)_u
                       -q_u(\sigma_u+1)O_0^2.
\]
The reflected expression uses $\sigma+5$ and $\sigma-1$.
Positivity of these gates implies $d<1/2$ and licenses that shortcut
only on its smaller domain. A bisection discards an interval only when
its entire defect lower bound exceeds the already acquired upper bound.

\subsubsection{Second moments and necessary noise floors}
The exact second-moment relations are
\begin{equation}\label{j:finite:beta-noise}
 \beta=\frac{v(1-2p)+s(2\alpha-1)}{v+s},\qquad
 m=\frac{v}{\sqrt{v+s}}.
\end{equation}
The second expression increases with $v$ and decreases with $s$.
For $s>0$, the first can be solved for $\alpha$; the case $s=0$
is tested directly. If an acquired upper bound $\beta_u$ is below
$t_l=1-2p_u$, then
\[
 v_l(t_l-\beta_u)\le s(\beta_u+1-2\alpha_l).
\]
A nonpositive right coefficient excludes the box; a positive one
gives a lower bound for $s$. If an acquired lower bound $\beta_l$
exceeds $t_u=1-2p_l$, use instead
$v_l(\beta_l-t_u)\le s(2\alpha_u-1-\beta_l)$.
These tests retain the same $\alpha$ as the conditional probabilities.

There is a further floor from the absolute third moment. The elementary
bounds $-x\le O_0s_0$, $y\le O_1s_1$ in
\eqref{j:core:conditional-moments} give
\[
 \rho_l(v+s)^{3/2}\le v(1-2v)+Ms,
\]
where $M$ bounds
$(3p+O_0)(1-\alpha)+(3q+O_1)\alpha$ at all four
$p,\alpha$ corners. Convexity of $(v+s)^{3/2}$ implies
\[
 A(v):=\rho_l v^{3/2}-v+2v^2
 \le [M-\tfrac32\rho_l\sqrt v]s.
\]
If $A(v_l)>0$, then
$\rho_l\sqrt{v_l}>1-2v_l$, so
$A'(v_l)>1/2+v_l>0$; moreover $A''>0$.
Hence a positive lower bound for $A(v_l)$ and an upper bound for
the right coefficient give a necessary noise floor on the whole box.
A nonpositive upper coefficient proves infeasibility. Directed square
roots preserve both conclusions.

Finally, the common inward bound $a_0+a_1\le d$, $d>0$,
together with \eqref{j:finite:conditional-moments}, implies
\begin{equation}\label{j:finite:common-inward}
 \rho(v+s)^{3/2}\ge v(1-2v)+B_*s+H_*s^2/d,
\end{equation}
where
\[
 B_*=3p(1-\alpha)+3q\alpha-d\max(\alpha,1-\alpha),\qquad
 H_*=\left(\frac{1-\alpha}{\sqrt q}+\frac\alpha{\sqrt p}\right)^2.
\]
For the quadratic term use
\[
 \frac{(1-\alpha)^2}{qa_0}+\frac{\alpha^2}{pa_1}
 \ge\frac{((1-\alpha)/\sqrt q+\alpha/\sqrt p)^2}{a_0+a_1}.
\]
The zero-variance restrictions cover vanishing endpoints by a limit.
The minimum of $B_*$ is at the four $p,\alpha$ corners, because it is
affine in $p$ and concave in $\alpha$.
On a $p$ interval $[p_l,p_u]$ a positive lower envelope for
$\sqrt{H_*}$ is $(1-\alpha)/\sqrt{q_u}+\alpha/\sqrt{p_u}$,
where $q_u=1-p_l$. This affine function is minimized at an
$\alpha$ endpoint. The global upper Taylor bound
\[
 (v+s)^{3/2}\le v^{3/2}+\tfrac32\sqrt v\,s
                                      +\tfrac38s^2/\sqrt v
\]
turns \eqref{j:finite:common-inward}, with $\rho\le\rho_u$, into a
necessary quadratic inequality in $s$. Its exact minimum on the
closed variance interval is tested at both endpoints and any interior
vertex. Strict positivity excludes the whole box. The implementation
defers $d=0$ to its deterministic branch.

\subsubsection{An additional outgoing event and full phase coverage}
For centered $W\in[-a_1,O_1]$ and $-a_1\le b<O_1$,
\[
 (W+a_1)(W-b)\le\ind_{\{W>b\}}(O_1+a_1)(O_1-b).
\]
Thus $\Pp(W>b)$ is at least the nonnegative part of
$(s_1-a_1b)/((O_1+a_1)(O_1-b))$.
For $b<-a_1$ the strict event is certain; for $b\ge O_1$ use zero.
If any of the $k$ class-one variables exceeds
$b=u+B+(k-1)a_1$, their sum exceeds $u+B$.
Independence gives the lower probability $1-(1-r)^k$, where $r$
is the preceding bound. The remaining centered sum, of variance $V_0$
and upper support $(n-k)a_0$, is at least $-B$ with probability at
least
\[
 \max\left\{\frac{B^2}{V_0+B^2},
                         \frac{B}{B+(n-k)a_0}\right\}\quad(B>0).
\]
These follow from Cantelli and Markov, respectively. At $B=0$ use
one if $V_0=0$ and zero otherwise. Multiplication of the two lower
bounds is valid because the two conditional groups are independent.
This supplies the optional single-event outgoing witness.

The acquired lower bound $m_l>0$ also gives
\[
 \lambda=n(v^2/m^2-v)\le n(v^2/m_l^2-v),\qquad h\le v_u/m_l.
\]
The quadratic on the right is convex, so its maximum on the $v$
interval occurs at an endpoint. If $-H_-\le H\le H_+$, the gates
\[
 np_l-H_-h_u>-1/2,\qquad np_u+H_+h_u<n+1/2
\]
ensure that all nearest counts belong to $\{0,\ldots,n\}$.
The original joint budget $d/d_0+\lambda/\Lambda\le1$ retains
its own intercept $\Lambda$ when a separate sharper noise cap clips
the count domain. Replacing that intercept by the sharper cap would
assert an additional joint inequality and is not licensed here.

\subsection{The signed fifth-index corner}
\label{j:finite:gapfree5}
The last small moment rectangle can be treated directly by conditioning
on signs. This avoids contact and Fourier estimates in that rectangle.
Let $\beta=\E[X|X|]$, $p=\Pp(X>0)$, $q=1-p$,
$m=\E X_+=\E X_-$ and $a_\pm=(1\pm\beta)/2$.
Zero atoms belong to the nonpositive class. Signwise Cauchy gives
\[
 m\ge\frac{1+\beta^2}{2\rho},\qquad
 \frac{m^2}{a_+}\le p\le1-\frac{m^2}{a_-},\qquad m^2\le pq.
\]
Set $h=pq/m$, $Y=p+hX$, $B=\ind_{\{X>0\}}$, and $U=Y-B$.
Both conditional means of $U$ vanish, with variances
\[
 v_0=p^2(qa_-/m^2-1),\qquad v_1=q^2(pa_+/m^2-1).
\]
Given $\sum_{i=1}^5B_i=j$, the noises are independent with variance
$V_j=(5-j)v_0+jv_1$. Write $b_j=\binom5j p^jq^{5-j}$ and define
\[
 T(V,d)=\begin{cases}1,&d\le0,\\V/(V+d^2),&d>0.\end{cases}
\]
Conditional Cantelli and the exact normal coordinate yield
\begin{equation}\label{j:finite:gapfree-bound}
 \frac{\sqrt5}{\rho}\{\Pp(S_5\le t)-\Phi(t/\sqrt5)\}
 \le\frac{\sqrt5}{\rho}
 \left\{\sum_{j=0}^5 b_jT(V_j,j-z)
          -\Phi\left(\frac{(z-5p)m}{\sqrt5pq}\right)\right\},
 \qquad z=5p+ht.
\end{equation}
The value $T(0,0)=1$ is essential: a deterministic sum at the threshold
belongs to the closed event.

\begin{proposition}[The positive corner]
\label{j:finite:gapfree5-result}
For every standardized iid marginal satisfying
\[
 \frac{477}{400}\le\rho\le\frac65,\qquad
 \frac{17}{50}\le\beta\le\frac{35139006743001}{10^{14}},
\]
the positive closed-CDF score in \eqref{j:finite:gapfree-bound} is strictly
below $4097/10000<\CE$ at every real threshold.
\end{proposition}
\begin{proof}
Here are the enclosures defining the finite proof. On a closed
$\beta$ slab $[b_l,b_u]$ put $m_l=(1+b_l^2)/(2(6/5))$.
The moment inequalities enclose $p$ in
$[2m_l^2/(1+b_u),\,1-2m_l^2/(1-b_l)]$.
On a closed $p,z$ box $[l,u]\times[a,b]\subset(0,1/2)\times[0,5]$
use
\begin{align*}
 v_0&\le\max\{0,u^2((1-l)(1-b_l)/(2m_l^2)-1)\},\\
 v_1&\le\max\{0,(1-l)^2(u(1+b_u)/(2m_l^2)-1)\}.
\end{align*}
Each binomial mass is maximized at the projection of $j/5$ onto
$[l,u]$, and each Cantelli factor is maximized at $z=b$ and the
variance upper bound. Clip the sum of the six upper bounds at one.
If $a-5u\ge0$, the normal argument is at least
$(a-5u)m_l/(\sqrt5\,u(1-u))$.
Otherwise it is at least $(a-5u)/\sqrt{5u(1-u)}$: use
$m\le\sqrt{pq}$ on the negative part and the fact that
$(z-5p)/\sqrt{pq}$ decreases in $p$ for $z\ge0$, $p<1/2$.
Its derivative has numerator $-[5p+z(1-2p)]/2$.
Outward square roots and a rational Gaussian lower enclosure therefore
bound the complete box. Take the nonnegative part of the CDF difference
before multiplying by $\sqrt5/\rho_l$.

The twelve closed $\beta$ slabs have endpoints
\[
 17/50,\quad 341/1000,\ldots,351/1000,\quad
 35139006743001/10^{14}.
\]
Together with the five intervals $z\in[k,k+1]$, $0\le k<5$, these
give sixty roots. The supplied binary subdivisions have 1,027 leaves.
Exact evaluation on every closed leaf gives a maximum upper bound
\[
 204848338979/500000000000<4097/10000.
\]
Both omitted rays have direct bounds. Use the full rectangle's $m_l$
and enclosing $[p_l,p_u]$, put $v_u=p_u(1-p_u)$ and
$s_u=v_u^2/m_l^2-v_u$. On the feasible domain $m^2\le v$,
$v^2/m^2-v$ increases with $v$ and decreases with $m$. Thus
$V_j\le\max\{j/p_l,(5-j)/(1-p_u)\}s_u$.
For $z\le0$, discard the nonnegative Gaussian CDF and evaluate the
Cantelli upper sum at zero, using binomial upper masses
$\binom5j p_u^j(1-p_l)^{5-j}$. The resulting score is below $331/1000$.
For $z\ge5$, the CDF is at most one and normal Cantelli gives the
score upper bound
\[
 \frac{\sqrt5}{\rho_l}\frac1{1+5m_l^2/p_u^2}<187/1000.
\]
These exact rational comparisons and the sixty roots cover all real $z$.
The proposition concerns the stated positive branch; it is not silently
extended to the opposite $\beta$ rectangle by reflection.
\end{proof}

\subsection{The numerical regions and the induction}
\label{j:finite:assembly}
The preceding lemmas reduce each finite index to a closed cover of
moment and probability parameters. Table~\ref{j:finite:region-table}
shows how the methods meet. All displayed intervals are closed; adjoining
rows overlap at their endpoints. The target of each Fourier region is
$T=4097/10000<\CE$. The affine bound closes the indicated upper tail;
at $\rho=7361/4000$ it is already at most $819/2000<T$.
\begin{table}[htbp]
\centering\small
\caption{The complete third-moment cover for indices five through thirty-nine.
The count column includes necessary-condition exclusions and, where specified,
the signed corner or finite-sine comparison.}
\label{j:finite:region-table}
\begin{tabular}{@{}cP{31mm}P{66mm}c@{}}
\toprule
$n$&Count region&Fourier complement and required support&Affine tail\\\midrule
5&$[1,6/5]$, with the positive corner
 &$[6/5,7361/4000]$; the six support rows of
 Proposition~\ref{j:finite:supportrows}&$\rho\ge7361/4000$\\
6--9&$[1,237/200]$
 &$[237/200,6/5]$, radius $13/5$; then the five support rows
 through $7361/4000$&$\rho\ge7361/4000$\\
10--14&$[1,23/20]$
 &$[23/20,71/50]$, radius $3$; $[71/50,15/8]$, unrestricted
 &$\rho\ge7361/4000$\\
15--19&$[1,23/20]$
 &$[23/20,33/25]$, radius $3$; $[33/25,19/10]$, unrestricted
 &$\rho\ge19/10$\\
20--24&$[1,9/8]$
 &$[9/8,6/5]$, radius $3$; $[6/5,19/10]$, unrestricted
 &$\rho\ge19/10$\\
25--34&$[1,6/5]$
 &$[6/5,19/10]$, unrestricted&$\rho\ge19/10$\\
35--39&$[1,6/5]$
 &$[6/5,19/10]$, unrestricted&$\rho\ge19/10$\\\bottomrule
\end{tabular}
\end{table}

The count regions split at $11/10$; at indices five through nine they
also split at $23/20$. At five the exact signed cover is as follows.
Put $b_2=84729828287999/(250000000000000)$ and
$b_3=35139006743001/10^{14}$. The signed-moment inequality gives
$|\beta|\le b_2$ up to $\rho=237/200$ and $|\beta|\le b_3$ up to
$\rho=6/5$: at each endpoint the rational test is
$(1+b_i^2)^2\ge\rho_i^2(1-b_i^2)$, with the difference increasing
in $b_i\ge0$. The count cover contains every feasible signed moment
for $1\le\rho\le477/400$, and every such moment with
$\beta\le17/50$ for $477/400\le\rho\le6/5$.
The remaining closed rectangle is exactly
Proposition~\ref{j:finite:gapfree5-result}; its positive sign is retained.
The splits at $237/200$, $467/400$, $477/400$, $1/5$ and $3/10$
are refinements of this union, not additional assumptions.

\paragraph{What the numerical assertion checks.}
The finite acquisition theorem~\ref{j:local:acquisition} transports
one actual marginal through every update. Support is proved before
higher moments are used, and an initial raw remainder before defect
feedback. Its simultaneous state includes the score, support, moments,
defect and signed threshold. A failed analytic gate gives no exclusion.
Each state must instead reach a necessary-condition contradiction, a
complete Fourier bound, or a complete count forest.

Write $P_j=[j/100,(j+1)/100]$, $0\le j<100$.
Every probability panel is closed. When several moment cells export
to it, upper bounds are maximized and lower bounds minimized over
\emph{all} those cells. Thus seams retain the source law. A count
forest covers every original root $k=0,\ldots,n$, including all signed
phase intervals. Missing panels are justified exclusions; refinements
must cover their entire parent. The exact partitions, score/HIGH caps,
refinement schedules and forest counts are certificate data. Their
logical content is the complete-domain assertion below, not the number
of successful files.

The following application-specific prerequisites are retained.
At five the first own-score bound is support-free; later Fourier inputs
use Proposition~\ref{j:finite:supportrows}, including every narrower
signed profile. For six through nine, the wider Fourier domain
$[23/20,6/5]$ uses radius $13/5$, while its low-band resource uses
radius $5/2$. At fifteen through nineteen, support radius three on
$\rho\le33/25$ precedes the own-score and HIGH bounds; at twenty,
radius three on $\rho\le6/5$ precedes them. At twenty-one through
twenty-four the own-score bounds are unrestricted. The twenty through
twenty-four initializers require no earlier numerical seed, and each
reused coarse count contract at twenty is admitted only after full
simultaneous containment has been established. A cellwise Fourier
bound evaluates the complete compensated expression, not a presumed
scalar HIGH ceiling. Each unrestricted region includes the complete
auxiliary-moment domain.

At $26$--$30$ and $32$--$34$, respectively, the acquired floors are
$N=25$ and $N=31$. Transfer retains $nD\le Nd$ as well as the raw
remainder bound. All binomial probabilities, normal arguments and
support tests still use the actual integer $n$. Local refinements of
a parent replace only the part they cover. At $26$--$30$ the surviving
second-band panels $P_{32},\ldots,P_{42}$ are closed by the finite-sine
inequality; the unrestricted Fourier complement makes a third count
band at twenty-five unnecessary. All omitted count panels, including
the two sign-mass exclusions at twenty and the weighted-budget
exclusions at $26$--$30$, have their own necessary-condition checks.

Here the five weighted exclusions have a direct explanation.
For $(n,j)=(30,52)$ and for $27\le n\le30$, $j=53$, the acquired
pairs $(d,\Theta)$ in $D\le25d/n$, $\theta\le\Theta$ are respectively
$(3/10,1/500)$ and $(5/16,1/2500)$. Put
$s=316227766/10^8<\sqrt{10}$ and $l=j/100$. Since
$p\ge l>p_*$, $\sqrt{p(1-p)}\le1/2$ and
$(1-\theta)^{3/2}\ge(1-\Theta)^2$, the joint budget gives
\[
 D\ge(12+4s)(1-\Theta)^2(l-(4-s)/2)^2
 >\begin{cases}251/1000,&j=52,\\38/125,&j=53.\end{cases}
\]
These contradict $D\le1/4$ at $(30,52)$ and
$D\le125/432<38/125$ at $j=53$, $n\ge27$.
The separate exclusions of $P_{36}$ and $P_{34}$ at twenty use
the sign-mass inequality on the respective bands, including endpoints.

\paragraph{The prerequisite at thirty-nine.}
Sharpness through 38 and the recurrence give
\[
 39(C_{39}-\CE)<55887/190000<3/10,\qquad C_{39}<.4226,
\]
with both deleted scores below $.4253$. Before the sharp-39 acquisition,
apply the separate comparison with actual $n=39$, $K=741/4000$,
predecessor allowance $.33$ and excess allowance $1/2$.
Under $C_{39}\ge\CE+K/39$, its support input is
$(N,L,K,q,R)=(39,39,741/4000,.33,3.84)$, and the positive charge
$K\rho_l$ remains in the resource. Its full acquisition excludes the
probability complement of $[.302,.458]$.
The joint fourth-moment test covers all phases for $0<\lambda\le.05$,
and nearest count covers $[.05,.08]$; on
$[.08,.10]$ use nearest count for $|u|\ge.19$ and whole count for
$|u|\le.19$; on $[.10,.12]$ use nearest count for $|u|\ge.27$
and whole count for $|u|\le.27$;
above $.12$ whole count covers every phase. Each band is clipped to its
own noise ceiling and includes its seams. Zero noise is Bernoulli.
The signed allowance tests of Section~\ref{j:local:allowance} exclude
the hypothesized score. Hence $C_{39}<.41449$, since
$\CE<.40974$ and $K/39=.00475$.

At 35--38, the entry bound
$C_n<41/100+4177/(13824n)$ supplies the own-score caps; at 39 use
the bound just proved. The initial support radii are $39/10$ at 35
and $19/5$ at 36--39. The independent support bound $96/25<39/10$
licenses the first, and all five initial and resource calculations
use the actual index. Any local support refinement is proved on its
whole source interval before export. Thus sharpness at 39 is never
used to establish its own entry cap.

\begin{proposition}[Complete numerical finite base through thirty-nine]
\label{j:finite:numerical539}
The support, Fourier, initialization, resource and count inequalities
in the preceding closed cover hold on their complete domains, including
the separate positive-39 prerequisite. Every exported row implies its
entire original count contract. The necessary exclusions and retained
forests cover every original moment and probability panel.
\end{proposition}
This assertion combines acquisition for one marginal, directed
containment of every scalar and signed-phase condition, count acceptance
at the actual integer, and the complete closed union. In particular,
if $q_i(H)$ are the acquired phase polynomials and
$\widetilde q_j(H)$ the contract polynomials, the join proves, for each
$i$, some $j$ with $q_i\le\widetilde q_j$ throughout that signed interval.
Endpoints and the possible interior minimum of the quadratic difference
suffice. Hence $\max_iq_i\le\max_j\widetilde q_j$. Selecting one row
retains the whole contract's actual index and moment domain.

\begin{proof}[Derivation of $C_n\le\CE$ for $5\le n\le39$]
Choose the least failing index and an attained positive closed maximizer.
Its predecessor is sharp, so contact loss is zero and every larger
retained allowance is valid. The table covers every $\rho\ge1$.
Fourier and signed-corner regions have score below $T<\CE$.
Otherwise acquisition retains the same law in an exported panel.
It cannot satisfy a necessary-condition exclusion; the directed join
therefore puts it in a complete count root, whose bound is below $\CE$.
At 39 the positive comparison is applied first. Every case contradicts
the assumed failure.
\end{proof}

\subsection{The actual indices forty, forty-one and forty-two}
\label{j:finite:4042}
A first failure here has sharp predecessors through two deletions.
The recurrence gives
\[
 0<n(C_n-\CE)<1433/4875<3/10,\qquad C_n<167/400<.4226.
\]
Thus the deleted bounds are below $.4253$, and the weaker allowances
$n(C_n-\CE)\le.54$ and contact loss $.19$ are valid.
The floor-40 argument preserves $\mathcal R_n\le r$ and $nD\le40d$;
its support test uses both endpoints 40 and 42 and proves radius $96/25$
before higher moments are used. Complete acquisition leaves left rows
covering $[.282,.38]$ and central rows covering $[.38,.476]$; their
complements are excluded. Confinement starts with widths
$.55,.575,.6$ on the left and $.55$ centrally, and cubic constants
30 and 24, with every width and positive-slope gate checked.

Both sides use the joint fourth-moment test for $0<\lambda\le.05$.
The central phase cover then uses nearest count on $[.05,.08]$,
then nearest count for $|u|\ge.19$ and whole count for $|u|\le.19$
on $[.08,.10]$, and nearest count for $|u|\ge.27$ and whole
count for $|u|\le.27$ on $[.10,.12]$;
whole count applies above $.12$.
The left cover uses nearest count from $.05$ through $.0925$, then nearest
count for $|u|\ge.16$ and whole count for $|u|\le.16$ on
$[.0925,.10]$, and nearest count for $|u|\ge.25$ and whole count
for $|u|\le.25$ on $[.10,.12]$; then nearest count for $u\le-.27$
and whole count for $u\ge-.27$ on $[.12,.135]$; above $.135$
use whole count. Each closed band is clipped to its acquired noise
ceiling. Zero noise is Bernoulli.

These three indices admit a positive finite-binomial reserve. For
$[a,b]\subset(0,1/2)$ put
$F_j(p)=\Pr(\operatorname{Bin}(n,p)\le j)$ and
$z_j(p)=(j-np)/\sqrt{np(1-p)}$. Both decrease with $p$, while
$\sqrt{np(1-p)}/(1-2p(1-p))$ increases. Therefore both sides of
every jump are bounded by
\begin{equation}\label{j:finite:binomialslack}
 \frac{\sqrt{nb(1-b)}}{1-2b(1-b)}
 \max\{F_j(a)-\Phi(z_j(b)),\Phi(z_j(a))-F_{j-1}(b),0\},
 \qquad F_{-1}=0.
\end{equation}
Between jumps the CDF is constant, so $j=0,\ldots,n$ suffices.
Exact endpoint probabilities and outward Gaussian bounds on a complete
closed subdivision of $[.28,.48]$ give a positive uniform lower bound
$\epsilon_p$ for $\CE-C^{\rm bin}_{n,p}$ on each parent, at the actual
$n=40,41,42$.

Retain both signs in the finite corrections:
$I/n\le\max(I/40,I/42)$ and
$\sqrt n\pi_kJ\le\max(b_0J,J)$ when $b_0\le\sqrt n\pi_k\le1$.
If $S$ bounds $\rho_B(C-C_B)/\lambda$ on a positive-noise band
$[a,b]$, it suffices that $S<\rho_{B,l}\epsilon_p/b$.
For $S<0$ this is immediate; otherwise use $\lambda\le b$ and
$\rho_B\ge\rho_{B,l}$. The scalar bounds are convex in positive
noise, being affine terms plus nonnegative reciprocals of positive
affine terms; their two endpoint tests cover the band. No assertion
for all $n\ge40$ follows from this finite calculation.

\begin{proposition}[Numerical floor-40 assertion]\label{j:finite:numerical4042}
The preceding acquisition, two-endpoint support, actual-index binomial
reserve and complete signed local covers satisfy their strict tests
at $n=40,41,42$. Each local row consumes its acquired simultaneous
profile. Consequently $C_{40},C_{41},C_{42}\le\CE$.
\end{proposition}
\begin{proof}
The base through 39 and a least failure supply all entry premises.
Acquisition and confinement preserve its witness. The local comparison
with the signed finite corrections excludes every positive-noise case;
the Bernoulli theorem excludes zero noise.
\end{proof}

\begin{theorem}[The finite base]\label{j:finite:base}
For every integer $1\le n\le42$ and every standardized iid marginal
with finite third absolute moment, $\sqrt n\,\Delta_n(P)\le\CE\rho(P)$.
\end{theorem}
\begin{proof}
The starting bounds and the four-summand comparison establish indices
one through four. Proposition~\ref{j:finite:numerical539} and its
same-law induction establish five through thirty-nine;
Proposition~\ref{j:finite:numerical4042} supplies the last three indices.
\end{proof}
